% Options for packages loaded elsewhere
\PassOptionsToPackage{unicode}{hyperref}
\PassOptionsToPackage{hyphens}{url}
%
\documentclass[
]{book}
\usepackage{amsmath,amssymb}
\usepackage{iftex}
\ifPDFTeX
  \usepackage[T1]{fontenc}
  \usepackage[utf8]{inputenc}
  \usepackage{textcomp} % provide euro and other symbols
\else % if luatex or xetex
  \usepackage{unicode-math} % this also loads fontspec
  \defaultfontfeatures{Scale=MatchLowercase}
  \defaultfontfeatures[\rmfamily]{Ligatures=TeX,Scale=1}
\fi
\usepackage{lmodern}
\ifPDFTeX\else
  % xetex/luatex font selection
\fi
% Use upquote if available, for straight quotes in verbatim environments
\IfFileExists{upquote.sty}{\usepackage{upquote}}{}
\IfFileExists{microtype.sty}{% use microtype if available
  \usepackage[]{microtype}
  \UseMicrotypeSet[protrusion]{basicmath} % disable protrusion for tt fonts
}{}
\makeatletter
\@ifundefined{KOMAClassName}{% if non-KOMA class
  \IfFileExists{parskip.sty}{%
    \usepackage{parskip}
  }{% else
    \setlength{\parindent}{0pt}
    \setlength{\parskip}{6pt plus 2pt minus 1pt}}
}{% if KOMA class
  \KOMAoptions{parskip=half}}
\makeatother
\usepackage{xcolor}
\usepackage{longtable,booktabs,array}
\usepackage{calc} % for calculating minipage widths
% Correct order of tables after \paragraph or \subparagraph
\usepackage{etoolbox}
\makeatletter
\patchcmd\longtable{\par}{\if@noskipsec\mbox{}\fi\par}{}{}
\makeatother
% Allow footnotes in longtable head/foot
\IfFileExists{footnotehyper.sty}{\usepackage{footnotehyper}}{\usepackage{footnote}}
\makesavenoteenv{longtable}
\usepackage{graphicx}
\makeatletter
\def\maxwidth{\ifdim\Gin@nat@width>\linewidth\linewidth\else\Gin@nat@width\fi}
\def\maxheight{\ifdim\Gin@nat@height>\textheight\textheight\else\Gin@nat@height\fi}
\makeatother
% Scale images if necessary, so that they will not overflow the page
% margins by default, and it is still possible to overwrite the defaults
% using explicit options in \includegraphics[width, height, ...]{}
\setkeys{Gin}{width=\maxwidth,height=\maxheight,keepaspectratio}
% Set default figure placement to htbp
\makeatletter
\def\fps@figure{htbp}
\makeatother
\setlength{\emergencystretch}{3em} % prevent overfull lines
\providecommand{\tightlist}{%
  \setlength{\itemsep}{0pt}\setlength{\parskip}{0pt}}
\setcounter{secnumdepth}{5}
\usepackage{booktabs}
\usepackage{longtable}
\usepackage{booktabs}
\usepackage{longtable}
\usepackage{array}
\usepackage{multirow}
\usepackage{wrapfig}
\usepackage{float}
\usepackage{colortbl}
\usepackage{pdflscape}
\usepackage{tabularx}
\usepackage{xltabular}
\usepackage{threeparttable}
\usepackage{threeparttablex}
\usepackage[normalem]{ulem}
\usepackage{makecell}
\usepackage{xcolor}
\ifLuaTeX
  \usepackage{selnolig}  % disable illegal ligatures
\fi
\usepackage[]{natbib}
\bibliographystyle{plainnat}
\usepackage{bookmark}
\IfFileExists{xurl.sty}{\usepackage{xurl}}{} % add URL line breaks if available
\urlstyle{same}
\hypersetup{
  pdftitle={Probabilités et Statistiques appliquées à l'Hydrologie},
  pdfauthor={Benjamin Renard},
  hidelinks,
  pdfcreator={LaTeX via pandoc}}

\title{Probabilités et Statistiques appliquées à l'Hydrologie}
\author{Benjamin Renard}
\date{2026-11-27}

\begin{document}
\maketitle

{
\setcounter{tocdepth}{1}
\tableofcontents
}
\chapter*{Préambule}\label{pruxe9ambule}
\addcontentsline{toc}{chapter}{Préambule}

\chapter{Données et statistiques descriptives}\label{donnuxe9es-et-statistiques-descriptives}

Les chroniques de débit issues de stations hydrométriques constitueront les données de base utilisées dans ce document. Dans un souci de simplicité, la plupart des exemples utilisent des débits journaliers, mais les méthodes statistiques présentées s'appliquent également à des débits à d'autres pas de temps (horaire, mensuel ou pas de temps variable). De manière plus générale, elles s'appliquent également à d'autres variables hydrométéorologiques (pluies, températures, vent, etc.) et sont également utilisées dans de nombreux autres domaines (sociologie, economie etc.).

La plupart des analyses statistiques ne s'effectuent pas directement sur la chronique initiale, mais plutôt sur des séries de \textbf{variables hydrologiques} qui sont extraites, par exemple, de la chronique de débits journaliers. Ces variables ont pour but de caractériser de manière quantitative le phénomène étudié (les crues, les étiages, la ressource). L'objectif de ce chapitre est tout d'abord de présenter quelques approches permettant de définir des variables hydrologiques en basses, moyennes et hautes eaux. Cette présentation ne sera pas exhaustive (les possibilités sont infinies\ldots), et cherchera plutôt à illustrer la grande diversité de variables que l'hydrologue peut être amené à étudier. Dans un second temps, nous présenterons quelques méthodes de statistiques descriptives qui visent à résumer l'information contenue dans les séries de variables.

\section{Définition de variables hydrologiques}\label{duxe9finition-de-variables-hydrologiques}

En général, une variable hydrologique est calculée en deux étapes~:

\begin{itemize}
\tightlist
\item
  \textbf{Filtrage} de la chronique initiale. Par exemple, il est fréquent de travailler à partir de la série des débits moyennés sur \(d\) jours.
\item
  \textbf{Extraction} de la variable sur la série filtrée. Par exemple, on pourra calculer le minimum annuel, le nombre de jours passés sous un seuil de bas débit, etc.
\end{itemize}

\subsection{Filtrage et autre pré-traitements}\label{filtrage-et-autre-pruxe9-traitements}

Le pas de temps de la chronique initiale n'est pas forcément adapté à l'étude du phénomène d'intérêt. Par exemple, pour les bassins à crue lente, le débit de pointe à un pas de temps fin n'est pas forcément la variable la plus pertinente. De même, si l'on s'intéresse aux écoulements liés à la fonte nivale, on pourra souhaiter «~gommer~» les variations rapides de débit dues à des faibles pluies se superposant au signal de fonte.

Ces considérations conduisent parfois à filtrer la chronique initiale avant le calcul des variables. Les quatre filtres présentés dans ce document sont les suivants~:

\begin{itemize}
\tightlist
\item
  Moyenne mobile calculée sur une fenêtre de \(d\) jours.
\item
  Minimum mobile calculé sur une fenêtre de \(d\) jours.
\item
  \emph{Base Flow Separation} (BFS), qui vise à séparer le débit de base du débit de ruissellement.
\item
  \emph{Sequent Peak Algorithm} (SPA), qui calcule un déficit de volume cumulé.
\end{itemize}

Le premier filtre est largement connu et utilisé en routine pour le calcul de débits caractéristiques. Les filtres suivants sont moins utilisés en pratique, mais sont présentés afin d'illustrer d'autres choix possibles pour le calcul de variables hydrologiques.

Dans le cas du filtrage par moyenne ou minimum mobile, la série filtrée \(S_t\) est obtenue à partir de la chronique journalière initiale \(q_t\) de la manière suivante~:

\begin{align} 
\text{Moyenne mobile sur } d \text{ jours : } S_t&=\frac{1}{d} \sum_{k=t-d/2}^{t+d/2}q_k  \\
\text{Minimum mobile sur } d \text{ jours : } S_t&=\min_{k \in [t-d/2;t+d/2]}{q_k} 
\label{eq:filtering}
\end{align}

Dans l'équation \eqref{eq:filtering}, la fenêtre de calcul est centrée sur le pas de temps courant \(t\). Il est fréquent (notamment en statistiques) de calculer la moyenne ou le minimum sur une fenêtre \([t-d ;t]\) précédant le pas de temps courant \(t\). Ceci ne modifie que marginalement la série filtrée. Ces deux filtres ont pour effet de lisser les hydrogrammes en limitant les variations de débit à l'échelle journalière (cf.~exemples en Figure \ref{fig:filtering}). Le débit moyen sur \(d\) jours est proportionnel à un volume, alors que le débit minimum sur \(d\) jours peut être interprété comme la valeur continûment dépassée sur \(d\) jours consécutifs. Ce dernier filtre est notamment utilisé pour l'étude des crues.

\begin{figure}

{\centering \includegraphics[width=1\linewidth]{_main_files/figure-latex/filtering-1} 

}

\caption{Exemples de filtres appliqués à une chronique de débits journaliers (l'Azergues à Lozanne).}\label{fig:filtering}
\end{figure}

Les filtres BFS et SPA sont moins couramment utilisés pour le calcul de variables hydrologiques, et sont donc présentés plus en détail ci-dessous.
L'objectif du BFS est d'estimer le débit de base dans l'hydrogramme des débits journaliers. Plusieurs algorithmes existent dans la littérature. Nous présentons ici celui présenté par \citet{Tallaksen2004}~:

\begin{itemize}
\tightlist
\item
  La série de débits journaliers est découpée en \(n\) blocs consécutifs (non-recouvrant) de \(d\) jours.
\item
  Les valeurs minimales dans chaque bloc sont calculées, et sont notées \(m_1,\ldots,m_n\)
\item
  Pour \(k = 2,\ldots,n-1\)~:

  \begin{itemize}
  \tightlist
  \item
    Si \(w \times m_k < min(m_{k-1}, m_{k+1})\), \(m_k\) est considéré comme un point pivot.
  \end{itemize}
\item
  La série filtrée est finalement obtenue en interpolant linéairement entre les points pivots.
\end{itemize}

L'algorithme BFS dépend de deux paramètres \(d\) et \(w\). \citet{Tallaksen2004} recommandent des valeurs typiques égales à \(d\) = 5 jours et \(w = 0.9\), qui donnent généralement des résultats très satisfaisants.

La Figure \ref{fig:BFS} présente un exemple d'application de ce filtre à une série de débits journaliers, pour un bassin à forte influence nivale (la Durance à l'Argentière-la-Bessée, 963 km\(^2\)). Ce filtre est particulièrement intéressant si l'on s'intéresse à la caractérisation de l'onde de fonte nivale, puisqu'il permet de gommer les variations de débit vraisemblablement dues à des pluies. Par exemple, le maximum annuel de la série filtrée par BSF est une variable plus robuste pour caractériser l'onde de fonte que le maximum annuel des débits bruts.

Signalons que le BFS est également à la base d'un indice couramment utilisé, nommé le \emph{Base Flow Index} (BFI), définit comme le ratio entre le débit de base (estimé par l'algorithme BFS) et le débit total. Le BFI mesure ainsi la part des apports stockés (e.g.~apports souterrains ou écoulements de fonte) par rapport aux débits issus du ruissellement rapide.

\begin{figure}

{\centering \includegraphics[width=1\linewidth]{_main_files/figure-latex/BFS-1} 

}

\caption{Application du filtre BFS pour un bassin à influence nivale}\label{fig:BFS}
\end{figure}

Le filtre SPA est quant à lui plutôt orienté vers l'étude des étiages. Selon \citet{Tallaksen2004}, cet algorithme a été initialement développé pour le dimensionnement des réservoirs, mais il peut aussi être appliqué pour cumuler le déficit de volume d'une série de débits journaliers par rapport à un seuil de bas débit (que l'on appellera ici \(q_0\)).

Soient \(q_t\) le débit journalier et \(q_0\) un seuil de bas débit. La série filtrée par SPA est définie par~:

\begin{equation} 
S_t=\max(0;S_{t-1}+q_0-q_t)
\label{eq:SPA}
\end{equation}

La Figure \ref{fig:SPA} illustre l'application du filtre SPA. Les périodes durant lesquelles la série filtrée est nulle correspondent aux périodes de débits soutenus. Inversement, les événements d'étiages correspondent à des valeurs de SPA positives. Ce filtre possède donc la particularité intéressante d'~«~inverser~» l'hydrogramme, les plus fortes valeurs correspondant aux étiages les plus sévères. L'inconvénient de ce filtre est sa sensibilité aux données manquantes, du fait de sa définition récursive.

\begin{figure}

{\centering \includegraphics[width=1\linewidth]{_main_files/figure-latex/SPA-1} 

}

\caption{Exemple d’application du filtre SPA sur une chronique de débits journaliers.}\label{fig:SPA}
\end{figure}

En plus de l'utilisation de filtres, signalons également deux prétraitements importants utilisés en pratique~:

\begin{itemize}
\tightlist
\item
  Définition de l'\textbf{année hydrologique}~: suivant le phénomène auquel on s'intéresse, l'utilisation de l'année civile n'est pas compatible avec la périodicité naturelle du phénomène. Pour un bassin donné, si l'on s'intéresse aux hautes eaux, il sera judicieux de faire débuter l'année hydrologique au moins le plus sec (i.e.~ayant la moyenne interannuelle la plus faible), et inversement pour les basses eaux. Ce choix dépend évidemment du régime hydrologique~: en France métropolitaine, l'année hydrologique de hautes eaux débutera ainsi en fin d'été pour un bassin pluvial, mais au milieu de l'hiver pour un régime nival.
\item
  \textbf{Saisonnalisation}~: il arrive fréquemment de restreindre l'analyse à une période particulière de l'année. Par exemple, si l'on s'intéresse à la fonte nivale, on pourra se contenter d'étudier les débits printaniers et estivaux.
\end{itemize}

Dans la suite de cette section, ces prétraitements seront omis par souci de simplicité~: nous utiliserons systématiquement le terme «~année~», étant entendu que l'année pourra être remplacée par «~l'année hydrologique~» ou «~la saison~» suivant le phénomène étudié.

\subsection{Variables pour les moyennes eaux}\label{variables-pour-les-moyennes-eaux}

La variable la plus utilisée pour l'étude des moyennes eaux est le \textbf{débit moyen annuel} (parfois appelé \textbf{module annuel})~: il s'agit tout simplement de la moyenne des débits observés en une année donnée.

Il est également possible d'étudier des variables issues de la \textbf{courbe des débits classés} annuelle. Cette courbe est obtenue en classant, pour une année donnée, les débits par ordre décroissant (Figure \ref{fig:FDC}). On peut alors lire sur la courbe la durée annuelle de dépassement d'un débit donné. Des variables descriptives des moyennes eaux peuvent être extraites de cette courbe en calculant, par exemple, le débit dépassé pendant la moitié de l'année (\textbf{débit annuel médian}), ou plus généralement pendant \(p\)\% de l'année (en choisissant une valeur de \(p\) pas trop éloignée de 50\% si l'on s'intéresse aux moyennes eaux).

\begin{figure}

{\centering \includegraphics[width=1\linewidth]{_main_files/figure-latex/FDC-1} 

}

\caption{Exemples de courbes annuelles des débits classés.}\label{fig:FDC}
\end{figure}

\subsection{Variables pour les hautes eaux}\label{sec:HE}

Nous allons nous intéresser aux deux principales techniques utilisées par les hydrologues~pour l'étude des crues: l'échantillonnage par valeurs maximales annuelles (MAXAN) et l'échantillonnage par valeurs supérieures à un seuil (SUPSEUIL).

L'échantillonnage MAXAN consiste à sélectionner chaque année le débit observé le plus fort. L'échantillonnage SUPSEUIL consiste quant à lui à choisir un seuil de haut débit, puis à sélectionner les pointes des évènements dépassant ce seuil. Dans la pratique, plutôt qu'un seuil, on se donne en général un nombre d'évènements à sélectionner par an (en moyenne), et par itérations successives, on calcule le seuil conduisant à cet objectif.

Chacune de ces méthodes présente des avantages et des inconvénients, notamment dans l'optique d'une analyse statistique. La mise en œuvre de l'approche MAXAN est très simple. En ne sélectionnant qu'un unique événement par an, on évite le risque de sélectionner deux valeurs de débits correspondant au même événement de crue (le risque de sélectionner, par exemple, un débit le 31 décembre 2000 et un autre le 1er janvier 2001 est limité par l'utilisation d'années hydrologiques). L'inconvénient est d'ignorer un certain nombre d'évènements lors des années où beaucoup de crues se sont produites (par exemple, année 2004-2005 dans la Figure \ref{fig:MAXANvsSUPSEUIL}), et inversement de prendre en compte des évènements peu importants lors des années peu actives (par exemple, année 2006-2007). L'homogénéité de l'échantillon n'est donc pas optimale.

\begin{figure}

{\centering \includegraphics[width=1\linewidth]{_main_files/figure-latex/MAXANvsSUPSEUIL-1} 

}

\caption{Illustration des échantillonnages MAXAN et SUPSEUIL pour l'Azergues à Lozanne.}\label{fig:MAXANvsSUPSEUIL}
\end{figure}

L'approche SUPSEUIL est plus difficile à mettre en œuvre. En effet, il faut ajouter des contraintes d'indépendance afin de ne pas échantillonner plusieurs fois le même événement hydrologique (cf.~exemple en \ref{fig:SUPSEUIL}). On impose en général une contrainte d'espacement temporel minimal entre deux pointes sélectionnées, ainsi qu'une contrainte de redescente vers un débit de base. Bien choisies, ces contraintes permettent de garantir l'indépendance de l'échantillon (cf. \citet{Lang1995} pour un état de l'art sur la technique d'échantillonnage SUPSEUIL). Néanmoins, l'approche SUPSEUIL possède plusieurs avantages par rapport à l'échantillonnage MAXAN~:

\begin{itemize}
\tightlist
\item
  Il est possible d'étoffer l'échantillon disponible pour l'analyse statistique en choisissant, en moyenne, plus d'un événement par an.
\item
  Cet échantillon sera également plus homogène que celui fournit par la méthode MAXAN~: ainsi, il sera possible de ne sélectionner aucune valeur pour les années peu actives (par exemple, année 2006-2007 dans la Figure \ref{fig:MAXANvsSUPSEUIL}), et inversement d'en sélectionner plusieurs pour les années ayant connu plusieurs événements importants (par exemple, année 2004-2005).
\item
  L'approche SUPSEUIL conduit à s'intéresser au processus d'occurrence des crues~: on peut par exemple définir les variables «~nombre d'événements sélectionnés chaque année~», ou «~durée entre deux événements successifs~».
\end{itemize}

\begin{figure}

{\centering \includegraphics[width=1\linewidth]{_main_files/figure-latex/SUPSEUIL-1} 

}

\caption{Illustration de la nécessité des contraintes d’indépendance pour l’échantillonnage SUPSEUIL.}\label{fig:SUPSEUIL}
\end{figure}

Signalons que les approches MAXAN et SUPSEUIL peuvent s'appliquer sur des séries préalablement filtrées. On peut notamment appliquer les filtres moyenne / minimum mobile sur une durée \(d\) (avec par exemple \(d\) = durée caractéristique de crue). De plus, il existe de très nombreuses autres variables visant à caractériser les crues. Citons notamment~:

\begin{itemize}
\tightlist
\item
  Variables issues de la courbe annuelle des débits classés, par exemple, débit dépassé 5\% de l'année.
\item
  Volume annuel cumulé au-dessus d'un seuil de haut débit.
\item
  Variables décrivant les hydrogrammes de crue, par exemple, temps de montée / descente, asymétrie, etc.
\end{itemize}

\subsection{Variables pour les basses eaux}\label{BE}

Les étiages sont des phénomènes physiques complexes qui ne peuvent pas être complètement décrits sur la base d'une unique valeur numérique. L'ouvrage de \citet{Tallaksen2004} fournit un panorama très complet des différents indices utilisables pour décrire les étiages. En général, les caractéristiques d'étiages~les plus couramment étudiées sont la durée, le déficit de volume, et le pic (i.e.~débits les plus bas).

\subsubsection{Débit minimum annuel (MINAN)}\label{duxe9bit-minimum-annuel-minan}

Il correspond au pic de l'étiage, et constitue l'équivalent de l'approche MAXAN pour les basses eaux. Le débit minimum annuel est rarement calculé directement sur la chronique journalière (ou à pas de temps plus fin) initiale. On préfère généralement calculer le débit mensuel minimum (noté \textbf{QMNA}, correspondant au mois calendaire ayant le plus faible débit moyen), ou appliquer préalablement un filtre «~moyenne mobile~» sur \(d\) jours. Les durées les plus fréquemment utilisées (en France) sont 3, 7 et 30 jours, conduisant aux variables \textbf{VCN3}, \textbf{VCN7} et \textbf{VCN30}.

\subsubsection{Variables relatives à un seuil de bas débit (INFSEUIL)}\label{variables-relatives-uxe0-un-seuil-de-bas-duxe9bit-infseuil}

De manière similaire à l'approche SUPSEUIL utilisée pour l'étude des crues, on peut identifier des événements d'étiage comme des périodes durant lesquelles le débit est resté inférieur à un seuil donné. Néanmoins, contrairement à l'approche SUPSEUIL, il est très difficile de sélectionner plusieurs étiages indépendants au cours d'une même année. La Figure \ref{fig:INFSEUIL} illustre la définition de ces événements d'étiage~: au cours d'une année hydrologique donnée, sept événements sont identifiés dans cet exemple. Il est néanmoins discutable d'étudier ces sept événements séparément~: d'un point de vue hydrologique, il s'agit en fait du même événement d'étiage, entrecoupé par quelques remontées de débit dues à des pluies isolées. Pour cette raison, on préférera cumuler les caractéristiques de ces sept événements sur l'année. Plus précisément, on peut calculer les variables suivantes:

\begin{itemize}
\tightlist
\item
  \textbf{Déficit de volume}~: Il s'agit du déficit de volume par rapport au seuil de bas débit, cumulé sur la totalité de l'année hydrologique (aire en rouge sur la Figure \ref{fig:INFSEUIL}).
\item
  \textbf{Durée de l'étiage}~: Elle est définie comme le nombre de jours où le débit est resté inférieur au seuil de bas débit (somme des durées des sept événements d'étiage sur la Figure \ref{fig:INFSEUIL}).
\end{itemize}

\begin{figure}

{\centering \includegraphics[width=1\linewidth]{_main_files/figure-latex/INFSEUIL-1} 

}

\caption{Identification des périodes d’étiages à partir d’un seuil de bas débit pour pour l’Azergues à Lozanne.}\label{fig:INFSEUIL}
\end{figure}

\subsubsection{Variables basées sur la série filtrée par SPA}\label{variables-basuxe9es-sur-la-suxe9rie-filtruxe9e-par-spa}

La série filtrée par SPA fournit un moyen naturel de définir des événements d'étiage et d'en déduire des variables caractéristiques. Par exemple, on pourra définir un événement d'étiage comme une période durant laquelle la série SPA est restée strictement positive. On peut alors aisément calculer la durée de l'événement, tandis que le maximum de la série SPA au cours de l'événement décrit le déficit de volume cumulé au pic d'étiage. Par rapport aux variables de durée et de déficit de volume présentées précédemment, l'avantage est que cette approche évite de sectionner un étiage qui s'étendrait sur plusieurs années, comme illustré en Figure \ref{fig:SPA} dans le cas de l'étiage de 1921.

Concluons cette section en remarquant que la liste de variables d'étiage dressée dans ce document est loin d'être exhaustive. En particulier, on pourra encore une fois utiliser la courbe annuelle des débits classés pour extraire des variables descriptives de l'étiage (par exemple, débit dépassé 90\% de l'année).

\subsection{Variables de saisonnalité}\label{variables-de-saisonnalituxe9}

Les variables présentées jusqu'ici sont des variables de quantité. On s'intéresse également parfois à la saisonnalité des événements hydrologiques. Par exemple, les dates de début et de fin d'étiage peuvent être importantes en termes de gestion de la ressource en eau.

Des variables de saisonnalité peuvent aisément être déduites des variables de quantité exposées précédemment, en utilisant par exemple la date du maximum ou du minimum annuel, les dates de dépassement de seuil, etc. D'autres variables peuvent être calculées en utilisant le concept de \textbf{centre de masse}. Ce concept a été introduit par \citet{Stewart2005} dans le cadre de l'étude de la saisonnalité de l'onde de fonte pour des cours d'eau à influence nivale. Dans cette publication, le centre de masse est défini comme la date à laquelle 50\% du volume de fonte a été écoulé.

Cette définition peut aisément être généralisée d'une part pour l'étude d'autres types d'événements hydrologiques (crues ou étiages), d'autre part pour caractériser le début et la fin des événements (sans se limiter au «~pic~» de l'événement).

Dans le cas des étiages par exemple, on définira le centre de masse à \(p\)\% comme la date à laquelle \(p\)\% du déficit de volume annuel est atteint. En utilisant par exemple les centres de masse à 10\%, 50\% et 90\%, on définit respectivement le début, le centre et la fin de la période d'étiage (Figure \ref{fig:CoM}). On procèdera de même pour les crues, mais en considérant plutôt le volume écoulé au-dessus d'un seuil de haut débit.

\begin{figure}

{\centering \includegraphics[width=1\linewidth]{_main_files/figure-latex/CoM-1} 

}

\caption{Exemple d’indices de saisonnalité : dates de début, centre et fin d’étiage.}\label{fig:CoM}
\end{figure}

\subsection{Discussion~: choix des variables hydrologiques}\label{discussion-choix-des-variables-hydrologiques}

Dans certains cas le choix de la variable hydrologique à étudier est imposé. C'est le cas notamment pour les variables «~réglementaires~» (module ou QMNA pour les débits réservés ou les autorisations de prélèvement).

Si ce n'est pas le cas, le choix de la variable d'étude n'est ni anodin ni trivial. Par exemple, si l'on s'intéresse au régime des étiages, de nombreuses variables peuvent être calculées pour décrire la durée, l'intensité ou le déficit de volume des étiages (cf.~section \ref{BE}). Ainsi, ce que l'on appellera (de manière abusive car ambiguë) «~l'étiage décennal~», par exemple, correspondra à des phénomènes potentiellement bien différents suivant que l'on parle de durée, de volume ou d'intensité. En conséquence, il peut s'avérer intéressant de décrire le phénomène que l'on souhaite caractériser (crue ou étiages) par plusieurs variables hydrologiques.

Les principaux éléments guidant le choix du ou des variables à étudier sont les suivants~:

\begin{itemize}
\tightlist
\item
  \emph{Objectifs de l'analyse}. C'est certainement le premier aspect à évaluer~: étant donnés les enjeux, quel aspect du phénomène hydrologique cherche-t-on à caractériser en priorité~? Par exemple, certains ouvrages seront vulnérables à des débits de pointe élevés, d'autre à des durées de submersion élevées.
\item
  \emph{Propriétés du bassin versant à étudier}. A titre d'illustration, un débit de pointe horaire est pertinent pour l'étude des crues d'un petit bassin cévenol, mais est moins adapté à l'étude des bassins à crues lentes comme la Somme par exemple, pour lesquels on cherchera plutôt à caractériser les crues en termes de volume ou de durée.
\item
  \emph{Eviter la redondance d'information}. D'un point de vue plus statistique, il n'est pas forcément utile de sélectionner plusieurs variables très corrélées. Ainsi, si l'on cherche à caractériser le phénomène d'étiage dans son ensemble, on cherchera à utiliser quelques variables qui présentent des niveaux de corrélation limités.
\end{itemize}

\section{Statistiques descriptives}\label{statistiques-descriptives}

L'objectif des \textbf{statistiques descriptives} est de résumer l'information portée par une série de données représentant une des variables hydrologiques présentées précédemment, sur la base de quelques graphiques ou grandeurs caractéristiques.

\subsection{Représentations graphiques}\label{sec:cfc}

Il existe une infinité de façons de représenter graphiquement un jeu de données~: diagrammes en bâtons, courbes, nuages de point\ldots{} Tout dépend de la structure des données, mais également de l'information que l'on souhaite transmettre. La première représentation graphique dans le cas d'une série temporelle consiste généralement à tracer les données chronologiquement, comme illustré en Figure \ref{fig:Ncrues} pour la variable hydrologique ``Nombre annuel de crues''. Cette représentation permet une première visualisation des données qui peut parfois permettre de detecter des problèmes comme des données aberrantes ou des changements marqués dans la série.

\begin{figure}

{\centering \includegraphics[width=1\linewidth]{_main_files/figure-latex/Ncrues-1} 

}

\caption{Nombre de crues se produisant chaque année hydrologique sur l'Azergues à Lozanne}\label{fig:Ncrues}
\end{figure}

En plus du tracé chronologique des données, deux représentations sont importantes pour les données quantitatives car elles peuvent être reliées `auxà des notions probabilistes dont nous parlerons prochainement : l'\textbf{histogramme} et la \textbf{courbe des fréquences cumulées}.

L'histogramme consiste à représenter la \textbf{fréquence} (ou parfois l'\textbf{effectif}) de chaque valeur que l'échantillon peut prendre. L'effectif d'une valeur donnée est le nombre de fois où cette valeur a été observée, la fréquence est égale à l'effectif divisé par la taille de la série analysée. Par exemple, pour les données de la Figure \ref{fig:Ncrues}, la valeur «~0~» est d'effectif 13 et de fréquence 13/56 \(\approx\) 0.23.
Il est également possible de cumuler ces fréquences par ordre croissant des valeurs possibles~: par exemple fréquence(0) \(\approx\) 0.23, fréquence(1) \(\approx\) 0.27, et donc fréquence cumulée(1) \(\approx\) 0.23 + 0.27=0.5. La Table \ref{tab:NcruesTable} montre le calcul de ces quantités pour toutes les données de la Figure \ref{fig:Ncrues}.

La notion de fréquence cumulée est extrêmement importante puisque nous verrons ultérieurement que les notions de quantiles et de période de retour y sont étroitement liées. En hydrologie, la fréquence cumulée est souvent appelée \textbf{fréquence au non-dépassement} : en effet, la fréquence cumulée d'une valeur est égale au pourcentage de données dans l'échantillon qui sont inférieures ou égales à cette valeur. Notons qu'il peut sembler arbitraire d'utiliser «~inférieur ou égal~» plutôt qu'~«~inférieur strict~» dans la définition de la fréquence cumulée~: nous reviendrons sur ce point ultérieurement.

\begin{table}

\caption{\label{tab:NcruesTable}Effectifs, fréquences et fréquences cumulées.}
\centering
\begin{tabular}[t]{rrrr}
\toprule
Valeur & Effectif & Frequence & Frequence\_cumulee\\
\midrule
0 & 13 & 0.2321429 & 0.2321429\\
1 & 15 & 0.2678571 & 0.5000000\\
2 & 15 & 0.2678571 & 0.7678571\\
3 & 5 & 0.0892857 & 0.8571429\\
4 & 3 & 0.0535714 & 0.9107143\\
\addlinespace
5 & 2 & 0.0357143 & 0.9464286\\
6 & 1 & 0.0178571 & 0.9642857\\
7 & 2 & 0.0357143 & 1.0000000\\
8 & 0 & 0.0000000 & 1.0000000\\
\bottomrule
\end{tabular}
\end{table}

Les chiffres calculés dans la Table \ref{tab:NcruesTable} peuvent être représentés graphiquement comme illustré en Figure \ref{fig:NcruesFig}: à gauche, l'histogramme consiste à représenter la fréquence (ou parfois l'effectif) associée à chaque valeur, tandis que le graphique de droite montre la courbe des fréquences cumulées.

\begin{figure}

{\centering \includegraphics[width=1.1\linewidth]{_main_files/figure-latex/NcruesFig-1} 

}

\caption{Histogramme des fréquences et courbe des fréquences cumulées.}\label{fig:NcruesFig}
\end{figure}

Ces premiers calculs et graphiques élémentaires permettent d'apporter des éléments de réponse à des questions qui sont régulièrement posées à l'hydrologue, par exemple~:

\begin{itemize}
\tightlist
\item
  «~\emph{Quelle est la probabilité de ne pas observer de crue au cours de l'année}~?~»~: la fréquence observée de la valeur 0 valant 0.23, on serait tenté de répondre que cette probabilité vaut 0.23.
\item
  «~\emph{Quel nombre annuel de crue a une probabilité 1/10 d'être dépassé}~?~»~: on peut répondre à cette question en recherchant la valeur dont la fréquence au dépassement vaut 0.1, ou de manière équivalente dont la fréquence au non-dépassement vaut 1-0.1 = 0.9. En observant la Table \ref{tab:NcruesTable}, on constate que la fréquence cumulée 0.9 est atteinte~entre les valeurs de 3 et 4 crues par an.
\item
  «~\emph{Quelle est la probabilité d'observer 8 crues ou plus au cours d'une même année}~?~»~: Au cours des 60 années d'observation, on n'a jamais observé plus de 7 crues. On serait donc tenté de répondre que la probabilité d'en observer 8 ou plus est nulle. Néanmoins, on peut légitimement douter de cette estimation~: un événement qui n'a jamais été observé est-il pour autant de probabilité nulle, c'est à dire impossible~?
\end{itemize}

Le dernier point soulevé ci-dessus illustre une distinction importante. La question posée est en effet exprimée en termes de \emph{probabilité}. Or, la réponse est donnée en termes de \emph{fréquence}. Il se trouve que \textbf{les deux notions ne sont pas identiques}, ce qui explique la gêne exprimée dans le dernier point. Un événement peut très bien n'avoir jamais été observé (fréquence nulle), mais être de probabilité non-nulle~! La notion de fréquence se rapporte au monde des observations (qui sont en nombre limité), tandis que la notion de probabilité se rapporte à un monde plus théorique, gouverné par la théorie des probabilités, et que nous aborderons au chapitre suivant.

L'exemple du nombre annuel de crues utilisé jusqu'ici est particulier car il utilise une variable discrète. Un second exemple correspondant au cas plus courant d'une variable hydrologique continue est à présent utilisé, en l'occurrence les débits journaliers maximum annuels pour l'Azergues à Lozanne. Les données sont représentées dans la Figure ci-dessous.

\begin{figure}

{\centering \includegraphics[width=1\linewidth]{_main_files/figure-latex/maxanQt-1} 

}

\caption{Série chronologique des débits journaliers maximum annuels pour l'Azergues à Lozanne}\label{fig:maxanQt}
\end{figure}

Pour une variable continue, chaque valeur n'apparaît généralement qu'une seule fois (à quelques ex-aquo éventuels près). Le calcul de l'effectif et de la fréquence associée à chaque valeur n'est donc pas intéressant. Par contre on peut toujours calculer la fréquence au non-dépassement associée à chaque valeur. Pour la \(i^{ème}\) valeur (par ordre croissant, appelée également valeur de rang \(i\)) cette fréquence est égale à \(i/n\), où \(n\) est la taille de l'échantillon (égale à 55 ici). Ce calcul est détaillé dans le tableau ci-dessous.

\begin{table}
\caption{\label{tab:maxanTable}Valeurs, rangs et fréquences cumulées.}

\centering
\begin{tabular}[t]{rrr}
\toprule
Valeur & Rang & Fcumul\\
\midrule
15.1 & 1 & 0.018\\
19.5 & 2 & 0.036\\
27.4 & 3 & 0.055\\
28.4 & 4 & 0.073\\
29.7 & 5 & 0.091\\
\addlinespace
37.0 & 6 & 0.109\\
39.1 & 7 & 0.127\\
41.6 & 8 & 0.145\\
42.4 & 9 & 0.164\\
42.9 & 10 & 0.182\\
\addlinespace
44.0 & 11 & 0.200\\
46.4 & 12 & 0.218\\
48.4 & 13 & 0.236\\
50.7 & 14 & 0.255\\
50.7 & 15 & 0.273\\
\addlinespace
54.4 & 16 & 0.291\\
55.0 & 17 & 0.309\\
55.1 & 18 & 0.327\\
58.0 & 19 & 0.345\\
\bottomrule
\end{tabular}
\centering
\begin{tabular}[t]{rrr}
\toprule
Valeur & Rang & Fcumul\\
\midrule
64.5 & 20 & 0.364\\
66.5 & 21 & 0.382\\
69.6 & 22 & 0.400\\
70.5 & 23 & 0.418\\
73.6 & 24 & 0.436\\
\addlinespace
74.6 & 25 & 0.455\\
75.5 & 26 & 0.473\\
79.7 & 27 & 0.491\\
81.0 & 28 & 0.509\\
82.5 & 29 & 0.527\\
\addlinespace
82.5 & 30 & 0.545\\
84.0 & 31 & 0.564\\
85.5 & 32 & 0.582\\
86.3 & 33 & 0.600\\
86.6 & 34 & 0.618\\
\addlinespace
87.0 & 35 & 0.636\\
89.0 & 36 & 0.655\\
93.5 & 37 & 0.673\\
95.6 & 38 & 0.691\\
\bottomrule
\end{tabular}
\centering
\begin{tabular}[t]{rrr}
\toprule
Valeur & Rang & Fcumul\\
\midrule
97.7 & 39 & 0.709\\
98.8 & 40 & 0.727\\
99.2 & 41 & 0.745\\
105.0 & 42 & 0.764\\
106.0 & 43 & 0.782\\
\addlinespace
109.0 & 44 & 0.800\\
122.0 & 45 & 0.818\\
122.0 & 46 & 0.836\\
131.0 & 47 & 0.855\\
139.0 & 48 & 0.873\\
\addlinespace
140.0 & 49 & 0.891\\
161.0 & 50 & 0.909\\
164.0 & 51 & 0.927\\
166.0 & 52 & 0.945\\
233.0 & 53 & 0.964\\
\addlinespace
244.0 & 54 & 0.982\\
266.0 & 55 & 1.000\\
\bottomrule
\end{tabular}
\end{table}

On peut alors tracer comme précédemment la courbe des fréquences cumulées repréentée en Figure \ref{fig:maxanFig}(b). Pour l'histogramme, on effectue un regroupement des valeurs en classes, ce qui revient à «~discrétiser~» la variable continue. Le choix des classes reste arbitraire, même si certaines certaines règles empiriques ont été proposées (par exemple \url{https://fr.wikipedia.org/wiki/Histogramme\#Construction}). L'histogramme des fréquences pour des classes d'étendue 50 \(m^3.s^{-1}\) est donné en Figure \ref{fig:maxanFig}(a).

\begin{figure}

{\centering \includegraphics[width=1\linewidth]{_main_files/figure-latex/maxanFig-1} 

}

\caption{Histogramme des fréquences et courbe des fréquences cumulées pour les débits maximum annuels de l'Azergues à Lozanne}\label{fig:maxanFig}
\end{figure}

Comme précédemment, ces calculs permettent d'apporter des éléments de réponse à des questions qui sont régulièrement posées à l'hydrologue :

\begin{itemize}
\tightlist
\item
  «~\emph{Quel débit de crue a une probabilité 1/10 d'être dépassé}~?~»~(débit décennal) : on recherche la valeur ayant une fréquence au non-dépassement égale à 1-1/10=0.9. On trouve environ 160 \(m^3.s^{-1}\), ce qui fournit une première estimation.
\item
  «~\emph{Quel débit de crue a une probabilité 1/100 d'être dépassé}~?~»~(débit centennal) : on recherche la valeur ayant une fréquence au non-dépassement égale à 1-1/100=0.99. On peut alors donner la fourchette 244-266 \(m^3.s^{-1}\).
\item
  «~\emph{Quelle est la probabilité d'observer une crue supérieure à 200 \(m^3.s^{-1}\)}~?~» : en se basant sur la Figure \ref{fig:maxanFig}(b), on trouve une fréquence au non-dépassement d'environ 0.95. La probabilité recherchée est donc de l'ordre de 0.05.
\item
  «~\emph{Quelle est la probabilité d'observer une crue supérieure à 300 \(m^3.s^{-1}\)}~?~» : on serait tenté de répondre «~zéro~» puisque la valeur observée maximale est de 266 \(m^3.s^{-1}\)\ldots{} mais comme expliqué précédemment, cette estimation n'est pas raisonnable, et illustre une des limitations de l'approche basée sur la seule description empirique des données~: il faudrait se doter d'un outil pour pouvoir \textbf{extrapoler} au-delà de la valeur la plus forte. Nous y reviendrons dans la suite de ce document.
\end{itemize}

\subsection{Indicateurs numériques}\label{indicateurs-numuxe9riques}

En plus des représentations graphiques présentées précédemment, on souhaite généralement résumer les données analysées sur la base de quelques indicateurs quantitatifs, renseignant sur l'ordre de grandeur des observations, leur variabilité, etc.

\subsubsection{Indicateurs de localisation}\label{indicateurs-de-localisation}

Nous nous intéressons dans un premier temps aux indicateurs dits de localisation. Soient \((x_1,\ldots,x_n)\) les données à analyser. Le premier indicateur calculé est généralement la \textbf{moyenne}, souvent notée \(\bar{x}\).

\begin{equation} 
\bar{x}=\frac{1}{n}\sum_{i=1}^{n}x_i
\label{eq:mean}
\end{equation}

La \textbf{médiane} est la valeur qui sépare l'échantillon ordonné en deux sous-parties de même effectif. Par exemple, la médiane des valeurs 10, 15, 18, 19, 20 vaut 18. Si le nombre de valeurs est pair, on calculera la demi-somme entre les deux valeurs du milieu, soit une médiane de 16.5 pour l'échantillon des quatre premières valeurs ci-dessus. La médiane peut également être lue sur la courbe des fréquences cumulées~: c'est la valeur correspondant à une fréquence cumulée de 0.5. Par rapport à la moyenne, la médiane est beaucoup moins sensible à la présence de valeurs extrêmes dans l'échantillon. Par exemple, si l'on remplace la valeur «~20~» par «~2000~» dans l'exemple ci-dessus, la médiane est inchangée, alors que la moyenne sera évidemment très différente. L'exemple classique pour illustrer la différence entre la moyenne et la médiane est le salaire des Français~: la moyenne vaut environ 2730 €, la médiane 2180 € (salaires nets en 2024 dans le secteur privé, source INSEE).

Un troisième indicateur est parfois utilisé, surtout pour les données discrètes, il s'agit du \textbf{mode}, qui est la valeur (pas forcément unique) la plus fréquente dans l'échantillon. Pour des données quantitatives continues, il faut (comme pour l'histogramme) procéder à un regroupement des individus. L'inconvénient est que le mode devient dépendant du regroupement arbitraire effectué.

Le dernier indicateur est particulièrement important dans le domaine de l'hydrologie~: le \textbf{quantile d'ordre p} est la valeur de l'échantillon dont la fréquence cumulée vaut p.~Le quantile peut donc être lu sur la courbe des fréquences cumulées, en faisant éventuellement une interpolation linéaire entre deux points. La médiane est ainsi le quantile d'ordre 0.5.

Enfin, on reporte fréquemment le \textbf{minimum} et le \textbf{maximum} d'un échantillon.

\subsubsection{Indicateurs de dispersion}\label{indicateurs-de-dispersion}

Les indicateurs de dispersion ont pour but de décrire la variabilité des données.
La \textbf{variance} permet de mesurer la façon dont les données se dispersent autour de la moyenne :

\begin{equation} 
V_x=\frac{1}{n}\sum_{i=1}^{n}(x_i-\bar{x})^2
\label{eq:var}
\end{equation}

On trouve parfois une définition alternative dans laquelle le \(\frac{1}{n}\) est remplacé par \(\frac{1}{n-1}\). Ces deux définitions deviennent équivalentes lorsque n est grand.

L'\textbf{écart-type} est défini comme la racine carrée de la variance, et a l'avantage de s'exprimer dans la même unité que les données.

\begin{equation} 
S_x=\sqrt{V_x}= \sqrt{\frac{1}{n}\sum_{i=1}^{n}(x_i-\bar{x})^2}
\label{eq:sdev}
\end{equation}

Le \textbf{coefficient de variation} est également utilisé pour comparer en relatif la variabilité de plusieurs séries de données dont les ordres de grandeurs ne sont pas comparables. Il est plutôt adapté aux données positives (les débits par exemple !) et n'est pas défini pour les données dont la moyenne est nulle. Il est généralement exprimé en pourcentage.

\begin{equation} 
CV_x=\frac{S_x}{\bar{x}}
\label{eq:cv}
\end{equation}

\subsubsection{Autres indicateurs}\label{autres-indicateurs}

Mentionnons finalement quelques indicateurs supplémentaires, qui n'ont pas de signification aussi claire que les indicateurs de localisation ou de dispersion, mais qui sont importants en Statistiques, notamment dans le cadre de la théorie de l'estimation.

Le \textbf{moment d'ordre k} est défini de la manière suivante~:

\begin{equation} 
m_k=\frac{1}{n}\sum_{i=1}^{n}x_i^k
\label{eq:moment}
\end{equation}

Pour \(k=1\), on reconnaît la moyenne. On définit de même le \textbf{moment centré d'ordre k}~:

\begin{equation} 
m^\prime_k=\frac{1}{n}\sum_{i=1}^{n}(x_i-\bar{x})^k
\label{eq:cmoment}
\end{equation}

Pour \(k=2\), on reconnaît la variance. Les moments centrés sont utilisés pour calculer deux indicateurs relatifs à la forme de l'histogramme~:

\begin{equation} 
\textbf{Asymétrie}: A_x=\frac{m^\prime_3}{S_x^3}
\label{eq:skew}
\end{equation}

\begin{equation} 
\textbf{Aplatissement} \text{ ou } \textbf{Kurtosis} : K_x=\frac{m^\prime_4}{S_x^4}
\label{eq:kurt}
\end{equation}

Comme son nom l'indique, l'asymétrie décrit le caractère symétrique de l'histogramme (l'asymétrie est nulle pour des données distribuées symétriquement). L'aplatissement décrit quant à lui le caractère plus ou moins «~pointu~» de l'histogramme.

En guise d'illustration, le Tableau ci-dessous fournit un résumé statistique pour les maxima annuels de l'Azergues à Lozanne étudiées précédemment. Fournir ce type de résumé lors de toute analyse statistique est une bonne pratique.

\begin{table}
\centering
\caption{\label{tab:summaryTable}Résumé statistique des valeurs maximales annuelles pour l’Azergues à Lozanne.}
\centering
\begin{tabular}[t]{lr}
\toprule
Moyenne & 87.59\\
Minimum & 15.10\\
Quantile 25\% & 50.70\\
Médiane & 81.00\\
Quantile 75\% & 102.10\\
\addlinespace
Maximum & 266.00\\
Ecart-type & 53.00\\
Coefficient de variation & 0.61\\
Variance & 2808.99\\
Asymétrie & 1.51\\
\addlinespace
Applatissement & 5.52\\
N & 55.00\\
\bottomrule
\end{tabular}
\end{table}

\chapter{Probabilités}\label{probabilituxe9s}

Le chapitre précédent a permis de présenter quelques méthodes permettant de décrire une série d'observations et d'en extraire de l'information. Ce chapitre tente d'aller au-delà d'une simple description empirique des données, en utilisant des \textbf{modèles probabilistes} pour décrire la \textbf{distribution} des observations.

La théorie des probabilités fournit un cadre général pour décrire les caractéristiques de variables «~génératrices~» de données. La puissance de cette théorie tient à sa généralité. En contrepartie, la définition même du terme «~probabilité~» fait appel à des notions mathématiques trop élaborées pour être présentées dans ce cours si l'on souhaite être totalement rigoureux du point de vue mathématique. Nous nous contenterons donc d'étudier une version «~allégée~» de la théorie probabiliste, qui sera néanmoins amplement suffisante en pratique.

Cette seconde section est organisée de la manière suivante. Dans un premier temps, les principales motivations conduisant à s'intéresser au cadre probabiliste seront présentées (section \ref{motivations}). Les notions de base seront ensuite définies (section \ref{sec:theorie-proba}), avant de présenter plus en détail quelques distributions fréquemment utilisées en hydrologie~: la loi de Gauss ou loi normale (section \ref{sec:gauss}), la loi log-normale (section \ref{sec:logN}), et les lois de valeurs extrêmes (sections \ref{sec:GEV}-\ref{sec:GPD}). Un résumé d'autres distributions usuelles est également fourni en Annexe \ref{annexe:distributions}.

\section{Motivations}\label{motivations}

L'aproche descriptive présentée dans le chapitre précédent permettait déjà d'apporter des éléments de réponse aux questions fréquemment posées à l'hydrologue (par exemple, «~\emph{Quelle valeur de débit a une probabilité 1\% d'être dépassée chaque année}~?~»~; «~\emph{Quelle est la probabilité de dépassement de la crue de décembre 2003}~?~»). On peut alors s'interroger sur l'intérêt de recourir à des concepts mathématiques supplémentaires. Nous expliquons ici pourquoi il n'est pas raisonnable de se contenter de cette approche descriptive.

\subsection{Quelle formule pour la fréquence au non-dépassement~?}\label{quelle-formule-pour-la-fruxe9quence-au-non-duxe9passement}

Nous avons déjà mentionné en section \ref{sec:cfc} le caractère arbitraire de la formule utilisée pour la fréquence au non-dépassement. Par définition, la fréquence au non-dépassement de la valeur de rang \(i\) est égale à \(i/n\), où \(n\) est le nombre d'observations. Ceci correspond à la définition suivante~: «~La fréquence au non-dépassement d'une valeur donnée est égale à la fréquence des observations inférieures ou égales à cette valeur~». Mais si l'on remplace «~inférieur ou égal~» par «~inférieur strict~», la formule devient alors \((i-1)/n\). En plus de ces deux formules naturelles, la littérature regorge d'autres formules pour le calcul des fréquences au non-dépassement. On peut citer par exemple la formule de Hazen~: \((i-0.5)/n\). Cette formule est en quelque sorte « au milieu~» des deux formules naturelles \(i/n\) et \((i-1)/n\) (et nous conseillons d'utiliser cette formule par défaut). Mais il en existe de nombreuses autres, par exemple \(i/(n+1)\) (Weibull), \((i-0.31)/(n+0.38)\) (Beard), \((i-0.44)/(n+0.12)\) (Gringorten), etc. Chacune de ces formules résulte de diverses considérations théoriques (non détaillées ici), mais dans tous les cas, il serait donc souhaitable de se donner un cadre méthodologique s'affranchissant de ce choix largement arbitraire.

\subsection{Interpolation et extrapolation}\label{interpolation-et-extrapolation}

Les exemples donnés en section \ref{sec:cfc} ont permis d'illustrer les limites d'une approche purement empirique pour l'\textbf{extrapolation}. Ainsi, si l'on se basait uniquement sur les fréquences au non-dépassement empiriques (avec la formule \(i/n\)), on conclurait que la plus forte valeur observée ne peut pas être dépassée~! L'utilisation d'une autre formule de fréquence empirique ne résoudrait que partiellement le problème, puisque l'extrapolation au-delà de la valeur la plus forte resterait problématique. Mentionnons également que l'\textbf{interpolation} entre deux fréquences empiriques par de segments de droite n'est guère satisfaisante.

\subsection{Variabilité d'échantillonnage et incertitudes}\label{variabilituxe9-duxe9chantillonnage-et-incertitudes}

La plus forte motivation conduisant à l'utilisation de modèles probabilistes est liée à la \textbf{variabilité d'échantillonnage}. Ce terme générique désigne le fait que les observations de débit dont nous disposons ne constituent généralement qu'un faible échantillon de l'ensemble des réalisations possibles de la rivière. En conséquence, les estimations seront dépendantes des données effectivement observées.
Cette variabilité d'échantillonnage peut être illustrée en considérant les 28 maxima annuels représentés en Figure \ref{fig:var-echantillonnage} (gauche). On constate que deux événements forts ont été observés lors de la première moitié de la période. La figure de droite illustre l'impact de ces deux valeurs sur l'allure des courbes de fréquence cumulée calculées sur deux périodes de 14 ans. Ainsi, si l'hydrologue doit répondre à la question «~quel débit de crue a une probabilité au non-dépassement de 0.9~?~», une approche purement empirique conduira à répondre «~environ 100 \(m^3.s^{-1}\)~» en utilisant les données de la seconde période, mais «~environ 250 \(m^3.s^{-1}\)~» sur la base de la première~!

Précisons que l'utilisation de modèles probabilistes ne permet pas de s'affranchir totalement de cette difficulté~: en effet, leur estimation est également soumise à la variabilité d'échantillonnage, mais dans une moindre mesure. De plus, l'approche probabiliste possède l'avantage majeur de permettre la quantification de cette variabilité, et de l'\textbf{incertitude} qui en résulte.

\begin{figure}

{\centering \includegraphics[width=1\linewidth]{_main_files/figure-latex/var-echantillonnage-1} 

}

\caption{Illustration du concept de variabilité d'échantillonnage : valeurs maximales annuelles pour l'Azergues à Lozanne sur 2 périodes de 14 ans (gauche), et courbes des fréquences cumulées correspondantes (droite).}\label{fig:var-echantillonnage}
\end{figure}

\section{Théorie des probabilités}\label{sec:theorie-proba}

Cette section présente les concepts issus de la théorie des probabilités qui sont utiles à l'hydrologue.

\subsection{Probabilité}\label{probabilituxe9}

En termes intuitifs, une \textbf{probabilité} est un instrument de mesure du hasard~: elle évalue le caractère probable (ou improbable) d'\textbf{événements} que l'on peut (ou que l'on choisit) de considérer comme \textbf{aléatoires}.

Une probabilité varie entre 0 et 1~: plus la probabilité est proche de 1, plus l'événement a des chances de se produire~; inversement, une probabilité proche de zéro dénote un événement ayant peu de chance de se produire. Un événement impossible a une probabilité nulle, un événement certain a une probabilité égale à 1.

La notion de probabilité peut également être formalisée de façon plus mathématique. Une \textbf{probabilité} sur un ensemble \(\Omega\) est une fonction \(\Pr:\wp(\Omega) \mapsto [0;1]\) qui vérifie les trois axiomes suivants~:

\begin{enumerate}
\def\labelenumi{\arabic{enumi}.}
\tightlist
\item
  \(\Pr(A) \geq 0, \forall A \subset \Omega\)
\item
  \(\Pr(\Omega)=1\)
\item
  Soit \((A_i)\) une suite de sous-ensembles de \(\Omega\) deux à deux disjoints, c'est à dire que \(A_i \cap A_j=\emptyset, \forall i \neq j\). Alors:
  \[\Pr \left( \bigcup_{i=1}^{+\infty} A_i\right)=\sum_{i=1}^{+\infty} \Pr \left( A_i\right)\]
\end{enumerate}

Dans la définition ci-dessus, \(\Omega\) représente l'ensemble des issues possibles de l'expérience aléatoire, il est appelé \textbf{univers} ou \textbf{ensemble fondamental}. Un sous-ensemble A de \(\Omega\) est appelé un \textbf{événement}. La notation \(\wp(\Omega)\) représente l'ensemble des parties de \(\Omega\), c'est à dire l'ensemble de tous les sous-ensembles possibles de \(\Omega\), ou encore l'ensemble des événements.

Les trois axiomes des probabilités conduisent aux propriétés suivantes~:

\begin{itemize}
\tightlist
\item
  \(\Pr(A)\in [0;1], \forall A \subset \Omega\)
\item
  \(\Pr(\emptyset)=0\)
\item
  \(A \subset B \implies \Pr(A) \leq \Pr(B)\)
\item
  \(\Pr(A^c)=1 - \Pr(A)\)
\item
  \(\Pr(A \cup B)=\Pr(A)+\Pr(B)-\Pr(A \cap B)\)
\end{itemize}

\emph{Remarque 1}~: la définition donnée ci-dessus est en fait abusive d'un point de vue strictement mathématique, car il peut exister des sous-ensembles de \(\Omega\) pour lesquels la probabilité n'est pas définie (sous-ensembles non-mesurables). Il faut alors définir en ensemble d'évènements mesurables, ce qui requiert des outils mathématiques trop sophistiqués pour être évoqués dans ce cours.

\emph{Remarque 2}~: Un événement impossible a une probabilité nulle, un événement certain a une probabilité égale à 1. La réciproque n'est cependant pas vrai du point de vue mathématique~: des événements de probabilité nulle ne sont pas nécessairement impossibles, dans le sens où des sous-ensembles non vides de \(\Omega\) peuvent être de probabilité nulle. Par exemple, la probabilité de «~tirer au hasard~» un nombre entier dans l'ensemble des nombres réels est nulle, bien que l'ensemble des nombres entiers ne soit pas vide.

\subsection{Variable aléatoire}\label{variable-aluxe9atoire}

En termes intuitifs, une variable aléatoire permet de décrire de manière quantitative l'ensemble des résultats possibles d'une expérience aléatoire. Supposons par exemple que nous lancions deux dés, et que nous nous intéressions à la somme des deux faces observées. La variable aléatoire «~\(X\)~:= Somme des deux dés~» prend ainsi des valeurs entières comprises entre 2 et 12. Nous nous intéresserons au calcul de probabilité d'événements du type «~la somme des deux dés vaut 4~» (noté \(\Pr(X = 4)\)) ou «~la somme des deux dés est strictement inférieure à 6~» (noté \(\Pr(X < 6)\)).

D'un point de vue plus hydrologique, une variable aléatoire \(X\) représentera tout simplement la variable que nous souhaitons étudier, par exemple «~Nombre de crues observées en une année~» ou «~débit maximum annuel~» (cf.~exemples dans le chapitre précédent). Il est par contre important de bien faire la différence entre la variable aléatoire \(X\) et les observations \((x_1,\ldots,x_n)\)~: la variable aléatoire \(X\) est un objet mathématique abstrait, que l'on considère comme le «~générateur~» des données observées \((x_1,\ldots,x_n)\). On dit d'ailleurs que \((x_1,\ldots,x_n)\) sont des \textbf{réalisations} de la variable aléatoire \(X\). Les notations utilisées dans ce document utilisent des lettres capitales pour les variables aléatoires, des lettres minuscules pour leurs réalisations.

En termes plus mathématiques, une \textbf{variable aléatoire réelle} est une fonction définie sur l'univers \(\Omega\) et à valeurs réelles, \(X:\Omega \mapsto(-\infty;+\infty)\). Dans l'exemple des dés utilisé ci-dessus, l'ensemble \(\Omega\) est égal à l'ensemble des couples \((a, b)\), où \(a\) et \(b\) sont des entiers entre 1 et 6. \(X\) sera la fonction somme~:

\begin{align} 
X:(1,\ldots,6) \times (1,\ldots,6) &\mapsto (2,\ldots,12) \\
(a,b) &\mapsto a+b
\end{align}

Ainsi, l'événement «~la somme des deux dés vaut 4~» correspond à l'ensemble des couples \((a, b)\) pour lesquels \(a+b=4\), soit \(X^{-1}(\{4\})=\{(1, 3), (3, 1), (2, 2)\}\). Notons que nous utiliserons l'abus de notation \(X=4\) au lieu de la notation \(X^{-1}(\{4\})\). Dans cet exemple précis, \(X\) ne prend pas ses valeurs dans l'ensemble de tous les nombres réels, mais seulement dans un sous-ensemble de l'ensemble des entiers~: il s'agit d'une \textbf{variable aléatoire discrète}.

\subsection{Distribution d'une variable aléatoire discrète}\label{distribution-dune-variable-aluxe9atoire-discruxe8te}

Soit \(X\) une variable aléatoire discrète, à valeurs dans un ensemble dénombrable ordonné \(\{x_i,i=1,2,\ldots  \}\). La \textbf{loi de probabilité} de \(X\) est la donnée des quantités \(f(x_i)=\Pr(X=x_i), \forall i=1,2,\ldots\)

La \textbf{fonction de répartition} de \(X\) est définie par~: \(F(x_i)=\Pr(X \leq x_i), \forall i=1,2,\ldots\)

Ces deux fonctions sont illustrées dans la Figure \ref{fig:pmf}.

\begin{figure}

{\centering \includegraphics[width=1\linewidth]{_main_files/figure-latex/pmf-1} 

}

\caption{Loi de probabilité et fonction de répartition pour une variable aléatoire discrète.}\label{fig:pmf}
\end{figure}

Il existe une analogie entre la loi de probabilité représentée en Figure \ref{fig:pmf} et l'histogramme présenté dans le chapitre précédent, où nous avions reporté les fréquences (en ordonnées) de chaque observation (en abscisses): nous aurons l'occasion de revenir à plusieurs reprises sur cette analogie fréquence / probabilité. L'objectif d'un modèle probabiliste est justement de remplacer les fréquences observées par des probabilités issues du modèle, pour les raisons exposées en section \ref{motivations}.

Concernant la fonction de répartition, il est aisé de démontrer la relation suivante~: \(F(x_i)=\sum_{k=1}^{i}f(x_i)\). Autrement dit, la fonction de répartition est égale au cumul des probabilités individuelles de chaque valeur possible. On retrouve encore une fois l'analogie avec la courbe des fréquences cumulées~: la fonction de répartition calculée en une valeur \(x\) est l'analogue de la fréquence au non-dépassement d'une observation \(x\).
En corollaire, \(F\) est une fonction croissante, qui part de \(0\) et tend vers \(1\) en l'infini.

De ces deux définitions découlent quelques grandeurs caractéristiques de la variable aléatoire \(X\) (notez encore une fois l'analogie avec le chapitre précédent !)~:

\begin{itemize}
\tightlist
\item
  L'\textbf{espérance}, \(E(X)=\sum_{i=1}^{+\infty}x_i f(x_i)\)
\item
  La \textbf{variance}, \(Var(X)=\sum_{i=1}^{+\infty}(x_i-E(X))^2 f(x_i)\), et l'\textbf{écart type}, \(\sigma(X)=\sqrt{\left( Var(X) \right)}\)
\item
  Le \textbf{quantile d'ordre \(p\)}, noté \(x_p\), qui vérifie \(F(x_p)=p\)
\item
  Les \textbf{moments d'ordre \(k\)}, \(\mu_k(X)=\sum_{i=1}^{+\infty}x_i^k f(x_i)\). Si \(k=1\), on reconnaît l'espérance.
\item
  Les \textbf{moments centrés d'ordre \(k\)}, \(\mu^\prime_k(X)=\sum_{i=1}^{+\infty}(x_i-E(X))^k f(x_i)\). Si \(k=2\), on reconnaît la variance.
\end{itemize}

\emph{Remarque 1}~: notez la distinction, en termes de vocabulaire, entre l'espérance (de la variable aléatoire) et la moyenne (des données). Malheureusement, cette distinction explicite s'arrête là~et le même mot variance (par exemple) est utilisé à la fois pour la variable aléatoire et pour les données. On utilise donc un adjectif pour marquer cette différence~: on parle de variance \textbf{théorique} ou «~\textbf{de la population}~» pour une variable aléatoire, et de variance \textbf{empirique}, \textbf{expérimentale} ou «~\textbf{de l'échantillon}~» pour celle calculée sur les données.

\emph{Remarque 2}~: Etant donné que la fonction de répartition \(F\) est en escalier, le quantile d'ordre \(p\) n'est pas forcément défini (si \(p\) tombe «~entre les marches~» de l'escalier).

\emph{Remarque 3}~: Les quantités faisant intervenir des sommes infinies peuvent ne pas exister.

\subsection{Distribution d'une variable aléatoire continue}\label{distribution-dune-variable-aluxe9atoire-continue}

Les notions de loi de probabilité et de fonction de répartition décrites précédemment peuvent se généraliser au cas d'une variable continue, avec quelques différences liées au passage du discret au continu. Nous allons commencer cette fois ci par définir la distribution d'une variable aléatoire continue X comme la donnée d'une \textbf{fonction de répartition} \(F\), telle que \(F(x)=\Pr(X \leq x)\). La \textbf{densité de probabilité} de \(X\) sera alors la fonction \(f\) telle que~: \(f(x)=\int_{-\infty}^{x}F(u)\text{d} u\). En d'autres termes, la densité \(f\) est la dérivée de la fonction de répartition \(F\) (Figure \ref{fig:pdf}).

\begin{figure}

{\centering \includegraphics[width=0.75\linewidth]{_main_files/figure-latex/pdf-1} 

}

\caption{Densité de probabilité (haut) et fonction de répartition (bas) pour une variable aléatoire continue.}\label{fig:pdf}
\end{figure}

\emph{Propriété}~: si \(f\) est continue, alors \(\Pr(a<X \leq b)=\int_{a}^{b} f(x) \text{d}x = F(b)-F(a)\).

Dans le cas continu, on voit ainsi apparaître une analogie entre la probabilité d'un événement et l'aire sous la courbe de la densité. Cette analogie implique en particulier que \(\Pr(X=a)=0, \forall a \in \mathbb{R}\). C'est la raison pour laquelle il n'est pas possible de définir la loi de probabilité directement à partir des probabilités des éléments de \(\Omega\), ce qui conduit à raisonner sur des intervalles. On remarque encore une fois l'analogie avec les histogrammes dans le cas continu du chapitre précédent, où nous étions obligés de procéder à des regroupements en classes.

Une densité de probabilité est nécessairement une fonction positive (ou nulle), dont l'aire totale sous la courbe vaut 1. Par contre, rien n'interdit que certaines valeurs soient supérieures à 1~: en effet, la fonction \(f\) ne représente pas une probabilité, mais une densité (d'où son nom~!). La fonction de répartition est quant à elle une fonction croissante entre 0 et 1.

A partir de la densité de probabilité, il est possible de définir les mêmes grandeurs caractéristiques que dans le cas discret, en remplaçant les sommes par des intégrales~:

\begin{itemize}
\tightlist
\item
  L'\textbf{espérance}, \(E(X)=\int_{-\infty}^{+\infty}x f(x)\text{d}x\)
\item
  La \textbf{variance}, \(Var(X)=\int_{-\infty}^{+\infty}(x-E(X))^2 f(x)\text{d}x\), et l'\textbf{écart type}, \(\sigma(X)=\sqrt{\left( Var(X) \right)}\)
\item
  Le \textbf{quantile d'ordre \(p\)}, noté \(x_p\), qui vérifie \(F(x_p)=p\). De manière équivalente, le quantile d'ordre \(p\) est égal à l'inverse de la fonction de répartition prise en \(p\), \(x_p=F^{-1}(p)\).
\item
  Les \textbf{moments d'ordre \(k\)}, \(\mu_k(X)=\int_{-\infty}^{+\infty}x^k f(x)\text{d}x\). Si \(k=1\), on reconnaît l'espérance.
\item
  Les \textbf{moments centrés d'ordre \(k\)}, \(\mu^\prime_k(X)=\int_{-\infty}^{+\infty}(x-E(X))^k f(x)\text{d}x\). Si \(k=2\), on reconnaît la variance.
\end{itemize}

\emph{Remarque}~: Encore une fois, rien ne garantit l'existence des intégrales infinies. De plus, certaines lois ne sont pas définies sur \(\mathbb{R}\) tout entier, ces intégrales doivent alors être réduites aux supports de ces lois.

\section{Distribution des valeurs «~moyennes~»~: la loi normale}\label{sec:gauss}

\subsection{Définition}\label{duxe9finition}

La \textbf{loi de Gauss} (également appelée \textbf{loi normale}) est centrale en statistiques, puisqu'un grand nombre de phénomènes ont un comportement approximativement Gaussien. De plus, un théorème que nous verrons ultérieurement affirme qu'une moyenne (ou une somme) de variables aléatoires tend (presque) toujours en distribution vers une loi normale. Cette loi dépend de deux paramètres \(\mu\) et \(\sigma\)~représentant la moyenne et l'écart-type de la distribution (Figure \ref{fig:gauss}). Elle est parfois utilisée en hydrologie pour modéliser des modules annuels.

La loi normale, notée \(\mathcal{N}(\mu,\sigma)\), possède les propriétés suivantes~:

\begin{align}
&f(x)=\frac{1}{\sqrt{2\pi}\sigma}\exp \left( -\frac{(x-\mu)^2}{2\sigma^2}\right) \\
&F(x)=\int_{-\infty}^{x}f(u)\text{d}u \text{ (pas d'expression analytique)}\\
&\sigma >0 \\
&E(X)=\mu\\
&Var(X)=\sigma^2
\label{eq:gauss}
\end{align}

\begin{figure}

{\centering \includegraphics[width=1\linewidth]{_main_files/figure-latex/gauss-1} 

}

\caption{Exemples de densités (gauche) et de fonctions de répartition (droite) pour la loi normale.}\label{fig:gauss}
\end{figure}

\subsection{Justification théorique}\label{justification-thuxe9orique}

Le rôle important joué par la loi de Gauss en statistiques appliquées peut s'expliquer par le théorème central limite. Ce théorème peut s'exprimer de la manière suivante. Soit \((X_1,\ldots,X_n)\) des variables aléatoires indépendantes et de même distribution (on notera \textbf{iid} pour \emph{indépendantes et identiquement distribuées}), d'espérance \(\mu\) et d'écart-type \(\sigma\). Alors, \textbf{quelle que soit la distribution des \(X_i\)} (pourvu qu'espérance et variance existent), la distribution de la statistique «~moyenne empirique~» \(\bar{X}=\frac{1}{n}\sum_{i=1}^n X_i\) tend vers une loi Gaussienne quand \(n\) augmente. En pratique, une trentaine d'observations (indépendantes) est souvent jugée suffisante pour appliquer l'approximation Gaussienne. Ce résultat très général implique que la distribution gaussienne est un choix naturel pour décrire des variables qui peuvent s'écrire comme la moyenne (ou la somme) d'un grand nombre de valeurs.

La Figure \ref{fig:tcl} illustre le théorème central limite. Le panneau de gauche montre l'histogramme de 10 000 valeurs simulées dans une loi uniforme : sa forme est très différente de celui d'une loi gaussienne. On s'intéresse à présent à la moyenne de 2 valeurs tirées dans une loi uniforme. Le panneau du milieu montre l'histogramme correspondant, et indique que déjà l'histogramme prend une forme en cloche qui se rapproche un peu plus de celui d'une loi gaussienne. Si l'on refait l'exercice avec la moyenne de 10 valeurs tirées dans une loi uniforme, l'histogramme ressemble très fortement à celui d'une loi gaussienne.

\begin{figure}

{\centering \includegraphics[width=1\linewidth]{_main_files/figure-latex/tcl-1} 

}

\caption{Illustration du théorème central limite.}\label{fig:tcl}
\end{figure}

\section{Distribution log-normale}\label{sec:logN}

La distribution \textbf{log-normale}, également appelée \textbf{loi de Galton}, est une distribution dérivée de la loi normale de la manière suivante~: la loi log-normale de paramètres \(\mu\) et \(\sigma\)~est la distribution de la variable \(\exp(X)\), où \(X\) suit une loi normale de paramètres \(\mu\) et \(\sigma\). Inversement, le logarithme népérien d'une variable suivant une loi log-normale suit une loi normale. De par sa définition, une variable suivant une loi log-normale est strictement positive. Cette distribution est utilisée en hydrologie notamment pour modéliser des variables d'étiage ou parfois des modules annuels.

La loi log-normale, notée \(\mathcal{LN}(\mu,\sigma)\), possède les propriétés suivantes~:

\begin{align}
&f(x)=\frac{1}{x\sqrt{2\pi}\sigma}\exp \left( -\frac{(\log(x)-\mu)^2}{2\sigma^2}\right) \text{ si } x > 0 \text{ (0 sinon)} \\
&F(x)=\int_{-\infty}^{x}f(u)\text{d}u \text{ (pas d'expression analytique)}\\
&x>0, \sigma >0 \\
&E(X)=\exp(\mu+\sigma^2/2)\\
&Var(X)=[\exp(\sigma^2)-1]\exp(2\mu+\sigma^2)
\label{eq:logN}
\end{align}

\begin{figure}

{\centering \includegraphics[width=1\linewidth]{_main_files/figure-latex/logN-1} 

}

\caption{Exemples de densités (gauche) et de fonctions de répartition (droite) pour la loi log-normale.}\label{fig:logN}
\end{figure}

\section{Distributions des valeurs extrêmes~: maxima}\label{sec:GEV}

Les distributions des valeurs extrêmes jouent un rôle central en hydrologie, notamment pour l'estimation des crues. Inversement, l'hydrologie a été une des motivations premières au développement d'une théorie des valeurs extrêmes, puisque les travaux précurseurs d'Emil Julius Gumbel avaient pour ambition de fournir un cadre théorique aux problèmes d'estimation des crues. Cette section focalise sur les distributions utilisées pour modéliser les maxima annuels (méthode d'échantillonnages MAXAN, section \ref{sec:HE}). Une autre section sera consacrée à l'étude des dépassements de seuil (méthode d'échantillonnage SUPSEUIL).

\subsection{Définitions}\label{duxe9finitions}

La \textbf{loi de Gumbel}, notée \(Gu(\mu,\lambda)\), est ~la première loi des valeurs extrêmes à avoir été utilisée pour l'étude des maxima annuels. Il s'agit d'une loi à deux paramètres, le \textbf{paramètre de position} \(\mu\) et le \textbf{paramètre d'échelle} \(\lambda\). Ce dernier paramètre est également appelé le \textbf{gradex} en hydrologie francophone. La loi de Gumbel possède les propriétés suivantes~:

\begin{align}
&f(x)=\frac{1}{\lambda}\exp \left( -\frac{x-\mu}{\lambda} -\exp \left( -\frac{x-\mu}{\lambda} \right) \right) \\
&F(x)=\exp \left( - \exp \left( -\frac{x-\mu}{\lambda} \right) \right)\\
&\lambda >0 \\
&E(X)=\mu + \gamma \lambda, \gamma \approx 0.5772 \text{ (constante d'Euler-Mascheroni)}\\
&Var(X)=\frac{\pi^2}{6}\lambda^2
\label{eq:gumbel}
\end{align}

La \textbf{loi généralisée des valeurs extrêmes}, notée \(GEV(\mu,\lambda,\xi)\),~généralise la loi de Gumbel~en lui ajoutant un troisième paramètre appelé le \textbf{paramètre de forme}. Elle possède les propriétés suivantes~:

\begin{align}
&\text{Soit } z=1-\frac{\xi(x-\mu)}{\lambda} \\
&f(x)= \frac{1}{\lambda} z^{\frac{1}{\xi}-1} \exp \left( - z^{\frac{1}{\xi}}\right) \text{ si } z > 0 \text{ (0 sinon)}\\
&F(x)=\exp \left( - z^{\frac{1}{\xi}}\right) \text{ si } z > 0 \\
&z>0,\lambda >0,\xi \neq 0 \\
&E(X)=\mu + \frac{\lambda}{\xi} \left( 1-\Gamma(\xi+1)\right) \text{ (existe seulement si } \xi> -1 \text{)}\\
&Var(X)= \left(\frac{\lambda}{\xi}\right)^2 \left( \Gamma(2\xi+1)-\Gamma^2(\xi+1)\right) \text{ (existe seulement si } \xi> -\frac{1}{2} \text{)}
\label{eq:GEV}
\end{align}

La Figure \ref{fig:gev} présente les densités et fonctions de répartitions de trois lois GEV, qui possèdent les mêmes paramètres de position et d'échelle mais diffèrent par leur paramètre de forme \(\xi\). Si \(\xi<0\), la loi GEV est dite «~de type Fréchet~»~: il s'agit d'une distribution bornée à gauche (i.e.~elle possède une borne inférieure), que l'on qualifie de «~distribution à queue lourde~» pour signifier que la densité chute lentement vers zéro en l'infini. Si \(\xi>0\), la loi GEV est dite «~de type Weibull~»~: elle est bornée à droite (i.e.~elle possède une borne supérieure). Enfin, lorsque \(\xi\) tend vers zéro, la densité de la loi GEV tend vers la densité d'une loi de Gumbel. On pourra donc considérer que la loi de Gumbel est un cas particulier de la loi GEV, obtenu en prenant un paramètre de forme nul. Notons qu'il s'agit là d'un léger abus de langage, car les formules de l'équation \eqref{eq:GEV} ci-dessus ne sont pas définies pour \(\xi=0\).

\begin{figure}

{\centering \includegraphics[width=1\linewidth]{_main_files/figure-latex/gev-1} 

}

\caption{Exemples de densités (gauche) et de fonctions de répartition (droite) pour la loi GEV.}\label{fig:gev}
\end{figure}

Le paramètre de forme joue un rôle majeur dans l'estimation de quantiles de crues, car il gouverne leur comportement asymptotique. A titre d'illustration, la ligne horizontale en Figure \ref{fig:gev} (droite) correspond à la probabilité 0.95. L'intersection avec les fonctions de répartition donne donc la valeur du quantile d'ordre 0.95~: suivant la valeur du paramètre de forme, ce quantile varie entre environ 140 (pour \(\xi=0.5\)), environ 175 (pour \(\xi=0\)), et jusqu'à plus de 250 (pour \(\xi=-0.5\)).

\subsection{Justification théorique}\label{justification-thuxe9orique-1}

De même que le rôle central de la loi normale pour les valeurs «~moyennes~» est dû à l'existence d'un théorème central limite, la loi GEV trouve sa justification dans un théorème limite appelé le théorème des valeurs extrêmes.

Le théorème des valeurs extrêmes est, en quelque sorte, l'équivalent du théorème central limite, mais où la statistique «~moyenne empirique~» \(\bar{X}_n=\frac{1}{n} \sum_{i=1}^{n}X_i\) est remplacée par la statistique «~maximum~» \(M_n=\max(X_1,\ldots,X_n)\). La formulation exacte du théorème est un peu délicate, nous le retiendrons donc sous la forme suivante~: si la distribution de \(M_n\) converge, alors c'est forcément vers une loi généralisée des valeurs extrêmes \(GEV(\mu,\lambda,\xi)\). Ce résultat est vrai quelle que soit la loi parente dont est issu l'échantillon \((X_1,\ldots,X_n)\) : ce théorème justifie donc l'utilisation de la loi GEV pour décrire le comportement probabiliste des maxima annuels.

La Figure \ref{fig:evt} permet d'illustrer cette convergence dans le cas d'une loi parente normale \(\mathcal{N}(0,1)\).
Cette figure a été obtenue de manière analogue à la Figure \ref{fig:tcl}. Le panneau de gauche montre l'histogramme de valeurs simulées dans la loi parente. Dans le panneau du milieu, on s'intéresse au maximum de 10 valeurs tirées dans la loi parente: l'histogramme correspondant indique que la distribution de ce maximum est asymétrique et n'est donc plus gaussienne. Si l'on refait l'exercice avec le maximum de 365 valeurs tirées dans la loi parente, on obtient un histogramme encore plus asymétrique et on peut montrer qu'il s'agit dans ce cas d'une distribution de Gumbel. Si l'on change de distribution parente, la distribution du maximum convergera toujours vers une loi GEV, mais par forcément avec un paramètre de forme nul (correspondant à une loi de Gumbel). Par exemple, une loi parente uniforme conduira à une loi GEV à paramètre de forme positif pour le maximum (distribution bornée à droite), tandis qu'une une loi parente de Student conduira à une loi GEV à paramètre de forme négatif (distribution à queue lourde).

\begin{figure}

{\centering \includegraphics[width=1\linewidth]{_main_files/figure-latex/evt-1} 

}

\caption{Illustration du théorème des valeurs extrêmes.}\label{fig:evt}
\end{figure}

\section{Distributions des valeurs extrêmes~: dépassements de seuil}\label{sec:GPD}

De même que la loi GEV et son cas particulier, la loi de Gumbel, sont des candidats naturels pour modéliser les maxima annuels, nous verrons dans cette section qu'il existe des distributions naturelles pour modéliser les dépassements de seuils obtenus par la méthode d'échantillonnage SUPSEUIL~: la loi de Pareto généralisée, et son cas particulier la loi exponentielle.

\subsection{Définitions}\label{duxe9finitions-1}

La \textbf{loi exponentielle}, notée \(Exp(x_0,\lambda)\), est une distribution à deux paramètres~: le \textbf{seuil} \(x_0\), et le \textbf{paramètre d'échelle} \(\lambda\)~(également appelé \textbf{gradex} en hydrologie francophone).

\begin{align}
&f(x)=\frac{1}{\lambda}\exp \left( -\frac{x-x_0}{\lambda} \right) \text{ si } x \geq x_0 \text{ (0 sinon)} \\
&F(x)=1-\exp \left( -\frac{x-x_0}{\lambda} \right) \text{ si } x \geq x_0 \text{ (0 sinon)} \\
&x \geq x_0,\lambda >0 \\
&E(X)=x_0 +  \lambda \\
&Var(X)=\lambda^2
\label{eq:expo}
\end{align}

La \textbf{loi de Pareto généralisée}, notée \(GPD(x_0,\lambda,\xi)\),~généralise la loi exponentielle~en lui ajoutant un troisième paramètre appelé le \textbf{paramètre de forme}. Elle possède les propriétés suivantes~:

\begin{align}
&\text{Soit } z=1-\frac{\xi(x-x_0)}{\lambda} \\
&f(x)= \frac{1}{\lambda} z^{\frac{1}{\xi}-1} \text{ si } x \geq x_0 \text{ et } z > 0  \text{ (0 sinon)}\\
&F(x)=1 - z^{\frac{1}{\xi}} \text{ si } x  \geq x_0 \text{ et } z > 0   \\
&x  \geq x_0, z>0,\lambda >0,\xi \neq 0 \\
&E(X)=x_0 + \frac{\lambda}{1+\xi} \text{ (existe seulement si } \xi> -1 \text{)}\\
&Var(X)= \frac{\lambda^2}{(1+\xi)^2(1+2\xi)} \text{ (existe seulement si } \xi> -\frac{1}{2} \text{)}
\label{eq:GPD}
\end{align}

La Figure \ref{fig:gpd} présent les densités et fonctions de répartition pour trois lois GPD. Le rôle du paramètre de forme est similaire à ce que l'on observait pour la loi GEV (section \ref{sec:GEV})~: lorsque \(\xi<0\), la loi GPD est bornée à gauche par le seuil et possède une queue lourde à droite. Lorsque \(\xi>0\), la loi GPD est bornée à gauche et à droite. Enfin, lorsque \(\xi\) tend vers zéro, la densité de la loi GPD tend vers la densité d'une loi exponentielle, ce qui conduit à considérer la loi exponentielle comme un cas particulier de la loi GPD, obtenu en prenant un paramètre de forme nul.

\begin{figure}

{\centering \includegraphics[width=1\linewidth]{_main_files/figure-latex/gpd-1} 

}

\caption{Exemples de densités (gauche) et de fonctions de répartition (droite) pour la loi GPD.}\label{fig:gpd}
\end{figure}

Par rapport à l'échantillonnage par la méthode MAXAN, l'échantillonnage SUPSEUIL présente l'avantage intéressant de pouvoir s'intéresser au processus d'occurrence des crues. Ainsi, si l'on considère la Figure \ref{fig:MAXANvsSUPSEUIL}), il est possible de dénombrer le nombre annuel d'événements de crue. Cette variable, que nous noterons \(N\), est en général modélisée à l'aide d'une \textbf{loi de Poisson} notée \(\mathcal{P}(\tau)\). Le paramètre \(\tau\) de la loi de Poisson est appelé le \textbf{taux d'occurrence}. La loi de Poisson possède les caractéristiques suivantes~:

\begin{align}
&\Pr(N=k)= \exp(-\tau) \frac{\tau^k}{k!}\\
&\Pr(N\leq k)= \exp(-\tau) \sum_{i=0}^{k}\frac{\tau^i}{i!}\\
&k  \text{ entier positif}, \tau >0\\
&E(N)=\tau\\
&Var(X)= \tau
\label{eq:Poisson}
\end{align}

La Figure \ref{fig:poisson} illustre la loi de probabilité et la fonction de répartition d'une loi de Poisson.

\begin{figure}

{\centering \includegraphics[width=1\linewidth]{_main_files/figure-latex/poisson-1} 

}

\caption{Exemples de lois de probabilté (gauche) et de fonctions de répartition (droite) pour la loi de Poisson}\label{fig:poisson}
\end{figure}

\subsection{Justification théorique}\label{justification-thuxe9orique-2}

L'utilisation du couple loi de Poisson + loi GPD pour la modélisation des dépassements de seuil (échantillonnage SUPSEUIL) est justifiée par une relation forte avec l'utilisation de la loi GEV pour les maxima (échantillonnage MAXAN). En effet, il est possible de démontrer que si le processus d'occurrence suit une loi de Poisson \(\mathcal{P}(\tau)\) et le processus de dépassement du seuil une loi \(GPD(x_0,\lambda_{GPD},\xi_{GPD})\) (échantillonnage SUPSEUIL), ces deux variables étant supposées indépendantes, alors le maximum annuel suit une loi \(GEV(\mu,\lambda_{GEV},\xi_{GEV})\) (MAXAN). De plus, les paramètres des distributions sont liés par les relations suivantes~:

\begin{align}
\xi_{GEV} &= \xi_{GPD} \\
\lambda_{GEV} &= \lambda_{GPD} \times \tau^{-\xi}\\
\mu &= x_0 + \frac{\lambda_{GPD}(1-\tau^{-\xi})}{\xi}
\label{eq:GEVGPD}
\end{align}

\chapter{Inférence}\label{infuxe9rence}

Ce chapitre présente quelques méthodes pour l'estimation des paramètres des distributions présentées dans le chapitre précédent. Un résumé des estimateurs classiques est également fourni en Annexe \ref{annexe:estimateurs}.

\section{Estimation de paramètres}\label{estimation-de-paramuxe8tres}

\subsection{Principes de l'estimation}\label{principes-de-lestimation}

\subsection{Formulaire des estimateurs usuels}\label{formulaire-des-estimateurs-usuels}

\subsection{Estimation par la méthode des moments}\label{estimation-par-la-muxe9thode-des-moments}

\subsection{Estimation par la méthode du maximum de vraisemblance}\label{estimation-par-la-muxe9thode-du-maximum-de-vraisemblance}

\subsection{Quel estimateur choisir~?}\label{quel-estimateur-choisir}

\section{Estimation de débits caractéristiques}\label{estimation-de-duxe9bits-caractuxe9ristiques}

\subsection{Quantiles et périodes de retour}\label{quantiles-et-puxe9riodes-de-retour}

\subsection{Calcul des débits caractéristiques}\label{calcul-des-duxe9bits-caractuxe9ristiques}

\subsection{Contrôle et Validation}\label{contruxf4le-et-validation}

\chapter{Incertitudes}\label{incertitudes}

\section{Les différentes sources d'incertitudes en hydrologie}\label{les-diffuxe9rentes-sources-dincertitudes-en-hydrologie}

\section{Quantification des incertitudes}\label{quantification-des-incertitudes}

\section{Quelques techniques de quantification probabiliste de l'incertitude d'échantillonnage}\label{quelques-techniques-de-quantification-probabiliste-de-lincertitude-duxe9chantillonnage}

\chapter{Tests}\label{tests}

\section{Principe général d'un test statistique}\label{principe-guxe9nuxe9ral-dun-test-statistique}

\section{Exemple détaillé}\label{exemple-duxe9tailluxe9}

\chapter{Annexes}\label{annexes}

\section{Distributions usuelles}\label{annexe:distributions}

\section{Estimateurs usuels}\label{annexe:estimateurs}

\section{Tests statistiques}\label{annexe:tests}

  \bibliography{book.bib,packages.bib}

\end{document}
